\chapter{Introduction}
\label{chap:introduction}
The conservation and effective management of wildlife populations depend crucially on accurate methods for monitoring animal abundance and diversity. Traditional approaches, including manual identification, tagging are often invasive, costly, time-consuming and impractical at large scales or in remote environments. In recent years, significant progress has been achieved through the application of automated computer vision methods for wildlife monitoring, particularly though individual animal re-identification (Re-ID). Animal Re-ID, which refers to the task of identifying and matching individual animals across multiple observations solely on visual appearance, has emerged as a viable non-invasive method for ecological studies and conservation efforts.\cite{re-id-overview} Despite the advancements, current visual Re-ID systems still face substantial challenges, especially when applied to real-world scenarios characterized by significant appearance variations, environmental clutter, occlusions, and limited training data. Current methods typically require extensive manual labeling effort or species-specific model training, restricting their applicability in general ecological monitoring and conservation task so developing an efficient, flexible Re-ID framework capable of generalizing to virtually any animal species without requiring significant manual annotation or retraining represents an important and currently unmet need in conservation technology. The main aim of this thesis is just that, to develop a generalizable, robust animal Re-ID and population estimation system capable of performing reliably across diverse animal species with minimal manual intervention. Specifically, this thesis attempts to create a method suitable for real-world scenarios characterized by limited labeled data, eliminating the requirement for extensive manual labeling which can often be very difficult and error-prone cause it requires a certain expertise. To achieve these goals, the research was structured around two interlinked objectives:

\begin{enumerate}
    \item \textbf{Evaluation and optimization of robust classification methodologies}: Before addressing the complex task of population estimation, it was essential to establish a strong foundation for animal Re-ID that generalizes effectively across species. Publicly available datasets, notably WildlifeReID-10k \cite{adam2025wildlifereid10k}, were leveraged for evaluation and testing. This phase involved extensive experimentation with robust feature extraction methods, including DISK \cite{tyszkiewicz2020disklearninglocalfeatures} and LightGlue \cite{lindenberger2023lightgluelocalfeaturematching} combined with dimensionality reduction and aggregation approaches such as PCA, GMM and Fisher Vectors.
    \item \textbf{Development of population size estimation method: } Following the establishment of a reliable and generalizable individual classification framework, the ultimate goal was population counting. Implementatin of the Nested Importance Sampling (NIS)\cite{perez2024humanintheloopvisualreidpopulation}, a human-in-the-loop statistical framework leveraging pairwise similarity scores produced by the general-purpose Re-ID model. Critically, this method requires minimal human input and does not require additional training or extensive labeling for new animal species or datasets, thus significantly reduicng manual labor and expertise requirements in practical scenarios. Several potential week points were discovered in the paper itself which were attempted to address in this thesis. 
\end{enumerate}

All experiments and analyses in this thesis are conducted on publicly available datasets, primarily through the WildlifeReID-10k benchmark \cite{adam2025wildlifereid10k}, in order to support reproducible evaluation and fair comparisons across methods.
